% ============================================================
% Question Bank: ch08-01 Populations and Samples
% Topic Title: Populations and Samples
% CCSS: 7.SP.A.1
% Questions: 27 (15 MC + 6 SA + 6 Graphical)
% ============================================================

% --- Multiple Choice Questions (15) ---

\begin{practiceQuestion}{8-1-q01}{mc}
\begin{questionText}
A veterinarian wants to know the average weight of dogs in a city. She weighs $60$ dogs that visit her clinic over one month. What is the population in this study?
\end{questionText}
\choiceA{The $60$ dogs weighed at the clinic}
\choiceB{All dogs in the city}
\choiceC{All pets in the city}
\choiceD{The veterinarian's clinic}
\correctAnswer{B}
\explanation{The population is the entire group of interest. Since the veterinarian wants to know about all dogs in the city, the population is all dogs in the city.}
\end{practiceQuestion}

\begin{practiceQuestion}{8-1-q02}{mc}
\begin{questionText}
A factory produces $5{,}000$ light bulbs per day. A quality inspector tests $75$ bulbs selected at random from the production line. Which term best describes the $75$ bulbs?
\end{questionText}
\choiceA{Population}
\choiceB{Census}
\choiceC{Sample}
\choiceD{Parameter}
\correctAnswer{C}
\explanation{A sample is a subset of the population selected for study. The $75$ bulbs are a sample of the $5{,}000$ produced.}
\end{practiceQuestion}

\begin{practiceQuestion}{8-1-q03}{mc}
\begin{questionText}
A student wants to survey the favorite pizza topping of all $650$ students at her middle school. She asks every $10$th student entering the building on Monday morning. What type of sampling method is this?
\end{questionText}
\choiceA{Convenience sampling}
\choiceB{Stratified sampling}
\choiceC{Systematic sampling}
\choiceD{Voluntary response sampling}
\correctAnswer{C}
\explanation{Systematic sampling selects members at regular intervals from a list or sequence. Choosing every $10$th student is a systematic method.}
\end{practiceQuestion}

\begin{practiceQuestion}{8-1-q04}{mc}
\begin{questionText}
A radio station asks listeners to call in and vote for their favorite song. Why might this sample be biased?
\end{questionText}
\choiceA{The sample is too large}
\choiceB{Only people who listen to that station and feel strongly will call}
\choiceC{The station plays too many songs}
\choiceD{Calling is a random process}
\correctAnswer{B}
\explanation{Voluntary response samples are biased because only people with strong opinions tend to participate. This does not represent all listeners equally.}
\end{practiceQuestion}

\begin{practiceQuestion}{8-1-q05}{mc}
\begin{questionText}
A nutritionist surveys $40$ people at a gym about their eating habits and concludes that most adults in the town eat healthy food. What is wrong with this inference?
\end{questionText}
\choiceA{The sample size is too small}
\choiceB{The sample is not representative of all adults in the town}
\choiceC{The nutritionist should have used a census}
\choiceD{There is nothing wrong with the inference}
\correctAnswer{B}
\explanation{People at a gym are likely more health-conscious than the general population. This is a convenience sample that is biased toward healthy eaters.}
\end{practiceQuestion}

\begin{practiceQuestion}{8-1-q06}{mc}
\begin{questionText}
A librarian wants to know what kinds of books middle school students prefer. She divides the school into grade levels and randomly selects students from each grade. What sampling method is this?
\end{questionText}
\choiceA{Convenience sampling}
\choiceB{Voluntary response sampling}
\choiceC{Systematic sampling}
\choiceD{Stratified sampling}
\correctAnswer{D}
\explanation{Stratified sampling divides the population into groups (strata) and then randomly selects from each group. Dividing by grade level and sampling from each grade is stratified sampling.}
\end{practiceQuestion}

\begin{practiceQuestion}{8-1-q07}{mc}
\begin{questionText}
Which of the following samples is most likely to be representative of all households in a neighborhood?
\end{questionText}
\choiceA{Surveying households that have dogs}
\choiceB{Surveying households near the school}
\choiceC{Assigning each household a number and using a random number generator to select $30$}
\choiceD{Surveying households with children under $12$}
\correctAnswer{C}
\explanation{A simple random sample gives every household an equal chance of being selected, making it the most representative method.}
\end{practiceQuestion}

\begin{practiceQuestion}{8-1-q08}{mc}
\begin{questionText}
A town has $3{,}200$ registered voters. A reporter surveys $80$ voters outside a grocery store. What fraction of the population was sampled?
\end{questionText}
\choiceA{$\frac{1}{40}$}
\choiceB{$\frac{1}{80}$}
\choiceC{$\frac{1}{32}$}
\choiceD{$\frac{1}{4}$}
\correctAnswer{A}
\explanation{Divide the sample size by the population: $\frac{80}{3{,}200} = \frac{1}{40}$.}
\end{practiceQuestion}

\begin{practiceQuestion}{8-1-q09}{mc}
\begin{questionText}
A researcher wants to study the sleep habits of teenagers. She places a survey link on a social media page popular with teens. Which type of bias is most likely present?
\end{questionText}
\choiceA{The sample is stratified}
\choiceB{Response bias from the survey questions}
\choiceC{Voluntary response bias — only teens who see and choose to respond will be counted}
\choiceD{The sampling is random}
\correctAnswer{C}
\explanation{Posting a survey online creates voluntary response bias because only those who see the post and choose to respond will be in the sample. This does not give every teenager an equal chance of being selected.}
\end{practiceQuestion}

\begin{practiceQuestion}{8-1-q10}{mc}
\begin{questionText}
A school has $450$ seventh graders and $400$ eighth graders. A counselor randomly selects $30$ seventh graders and $30$ eighth graders. Why might this sampling plan be questioned?
\end{questionText}
\choiceA{The total sample size is too large}
\choiceB{The number selected from each grade is not proportional to the grade sizes}
\choiceC{Random selection always produces bias}
\choiceD{The counselor should only survey one grade}
\correctAnswer{B}
\explanation{In a proportional stratified sample, each group should be represented in proportion to its share of the population. The seventh grade is larger but has the same sample size as the eighth grade.}
\end{practiceQuestion}

\begin{practiceQuestion}{8-1-q11}{mc}
\begin{questionText}
A convenience sample is one in which participants are chosen because they are:
\end{questionText}
\choiceA{The best representatives of the population}
\choiceB{Easy to reach or readily available}
\choiceC{Randomly selected from a list}
\choiceD{Chosen every $n$th person in a line}
\correctAnswer{B}
\explanation{A convenience sample selects individuals who are easiest to reach. This method is quick but often produces biased results because it does not give everyone in the population an equal chance.}
\end{practiceQuestion}

\begin{practiceQuestion}{8-1-q12}{mc}
\begin{questionText}
A manager wants to know if employees are satisfied with their work schedule. She surveys only part-time employees. The population she wants to study is all employees. What is the main problem?
\end{questionText}
\choiceA{The sample is too large}
\choiceB{The sample excludes full-time employees, making it unrepresentative}
\choiceC{She should use a voluntary response method}
\choiceD{Part-time employees always have the same opinions as full-time employees}
\correctAnswer{B}
\explanation{Excluding full-time employees means the sample does not represent the entire population. Part-time employees may have different opinions about scheduling.}
\end{practiceQuestion}

\begin{practiceQuestion}{8-1-q13}{mc}
\begin{questionText}
A marketing company mails a survey to $2{,}000$ randomly selected households. Only $250$ surveys are returned. The responses may be biased because:
\end{questionText}
\choiceA{The initial mailing was not random}
\choiceB{People who returned surveys may have different opinions from those who did not}
\choiceC{$250$ is always too many responses}
\choiceD{Mail surveys are the best method available}
\correctAnswer{B}
\explanation{Non-response bias occurs when a large portion of the sample does not respond. Those who reply may feel more strongly about the topic than those who do not.}
\end{practiceQuestion}

\begin{practiceQuestion}{8-1-q14}{mc}
\begin{questionText}
Which question below is most likely to produce biased survey results?
\end{questionText}
\choiceA{What is your favorite subject in school?}
\choiceB{How many hours do you sleep on a school night?}
\choiceC{Don't you agree that recess should be longer?}
\choiceD{Do you prefer cats or dogs?}
\correctAnswer{C}
\explanation{Leading questions push respondents toward a particular answer. The phrase ``Don't you agree'' suggests the expected response, introducing response bias.}
\end{practiceQuestion}

\begin{practiceQuestion}{8-1-q15}{mc}
\begin{questionText}
A teacher puts the names of all $120$ students in a hat and draws $20$. Which statement is true about this sample?
\end{questionText}
\choiceA{It is a voluntary response sample}
\choiceB{It is a convenience sample}
\choiceC{It is a simple random sample}
\choiceD{It is a stratified sample}
\correctAnswer{C}
\explanation{Drawing names from a hat gives every student an equal chance of being selected, which is a simple random sample.}
\end{practiceQuestion}

% --- Short Answer Questions (6) ---

\begin{practiceQuestion}{8-1-q16}{sa}
\begin{questionText}
A dentist surveys $35$ of her $700$ patients about flossing habits. What percentage of the population is in the sample?
\end{questionText}
\correctAnswer{$5\%$}
\explanation{Divide the sample size by the population: $\frac{35}{700} = 0.05 = 5\%$.}
\end{practiceQuestion}

\begin{practiceQuestion}{8-1-q17}{sa}
\begin{questionText}
A school has $900$ students. A researcher wants a sample of $\frac{1}{15}$ of the population. How many students should be in the sample?
\end{questionText}
\correctAnswer{$60$}
\explanation{Multiply the population by the fraction: $900 \times \frac{1}{15} = 60$ students.}
\end{practiceQuestion}

\begin{practiceQuestion}{8-1-q18}{sa}
\begin{questionText}
Name one advantage of using a stratified random sample instead of a simple random sample when studying students across different grade levels.
\end{questionText}
\correctAnswer{It ensures every grade level is represented.}
\explanation{Stratified sampling divides the population into groups (strata) and samples from each one. This guarantees that every subgroup — such as each grade level — is represented in the final sample.}
\end{practiceQuestion}

\begin{practiceQuestion}{8-1-q19}{sa}
\begin{questionText}
A store manager asks the first $25$ customers who arrive on Saturday morning about store hours. Identify the type of sampling used.
\end{questionText}
\correctAnswer{Convenience sampling}
\explanation{The manager chose participants because they were easy to reach (the first to arrive), not by a random process. This is convenience sampling and may not represent all customers.}
\end{practiceQuestion}

\begin{practiceQuestion}{8-1-q20}{sa}
\begin{questionText}
A park has $1{,}400$ annual visitors. If $84$ visitors are randomly surveyed, what fraction of the population is in the sample? Write your answer in simplest form.
\end{questionText}
\correctAnswer{$\frac{3}{50}$}
\explanation{Divide: $\frac{84}{1{,}400} = \frac{84}{1400}$. Simplify by dividing numerator and denominator by $28$: $\frac{3}{50}$.}
\end{practiceQuestion}

\begin{practiceQuestion}{8-1-q21}{sa}
\begin{questionText}
Explain why surveying students only during after-school clubs to learn about the hobbies of all students at a school could produce biased results.
\end{questionText}
\correctAnswer{Students in clubs may have different hobbies than students who are not in clubs.}
\explanation{Sampling only from after-school clubs misses students who do not participate in clubs. Those students may have different hobby preferences, making the sample unrepresentative of the whole school.}
\end{practiceQuestion}

% --- Graphical Questions (3 GMC + 3 GSA) ---

\begin{practiceQuestion}{8-1-q22}{gmc}
\begin{questionText}
A student surveyed classmates about how they get to school. The results are shown below.

\begin{center}
\begin{tabular}{|l|c|}
\hline
\textbf{Method} & \textbf{Number of Students} \\
\hline
Bus & $14$ \\
Car & $8$ \\
Walk & $6$ \\
Bike & $2$ \\
\hline
\textbf{Total} & $30$ \\
\hline
\end{tabular}
\end{center}

If the school has $600$ students, which prediction is best supported by the sample?
\end{questionText}
\choiceA{About $120$ students bike to school}
\choiceB{About $280$ students take the bus}
\choiceC{About $300$ students walk to school}
\choiceD{About $160$ students ride in a car}
\correctAnswer{B}
\explanation{In the sample, $\frac{14}{30}$ of students take the bus. Multiply: $\frac{14}{30} \times 600 = 280$. About $280$ students in the school are predicted to take the bus.}
\end{practiceQuestion}

\begin{practiceQuestion}{8-1-q23}{gmc}
\begin{questionText}
Four different sampling plans were used to survey students about lunch preferences. The diagrams below show which students were included in each sample (shaded).

\begin{center}
\begin{tikzpicture}[scale=0.45]
% Plan A – only top row
\node[above] at (2.5,4.5) {\textbf{Plan A}};
\draw (0,0) grid (5,4);
\foreach \x in {0,...,4} \fill[funBlue!50] (\x,3) rectangle +(1,1);
% Plan B – random
\begin{scope}[xshift=7cm]
\node[above] at (2.5,4.5) {\textbf{Plan B}};
\draw (0,0) grid (5,4);
\fill[funBlue!50] (1,3) rectangle +(1,1);
\fill[funBlue!50] (3,2) rectangle +(1,1);
\fill[funBlue!50] (0,1) rectangle +(1,1);
\fill[funBlue!50] (4,0) rectangle +(1,1);
\fill[funBlue!50] (2,1) rectangle +(1,1);
\end{scope}
% Plan C – left column
\begin{scope}[xshift=14cm]
\node[above] at (2.5,4.5) {\textbf{Plan C}};
\draw (0,0) grid (5,4);
\foreach \y in {0,...,3} \fill[funBlue!50] (0,\y) rectangle +(1,1);
\end{scope}
% Plan D – scattered
\begin{scope}[xshift=21cm]
\node[above] at (2.5,4.5) {\textbf{Plan D}};
\draw (0,0) grid (5,4);
\fill[funBlue!50] (0,3) rectangle +(1,1);
\fill[funBlue!50] (2,2) rectangle +(1,1);
\fill[funBlue!50] (4,1) rectangle +(1,1);
\fill[funBlue!50] (1,0) rectangle +(1,1);
\fill[funBlue!50] (3,3) rectangle +(1,1);
\end{scope}
\end{tikzpicture}
\end{center}

Which two plans appear to use a method that selects students from across the entire population?
\end{questionText}
\choiceA{Plans A and C}
\choiceB{Plans B and D}
\choiceC{Plans A and D}
\choiceD{Plans C and D}
\correctAnswer{B}
\explanation{Plans B and D select students scattered throughout the grid, representing different parts of the population. Plans A and C sample from only one row or column, which may not be representative.}
\end{practiceQuestion}

\begin{practiceQuestion}{8-1-q24}{gmc}
\begin{questionText}
A researcher took three random samples from a town and recorded how many people preferred tea over coffee.

\begin{center}
\begin{tabular}{|c|c|c|}
\hline
\textbf{Sample} & \textbf{Sample Size} & \textbf{Prefer Tea} \\
\hline
$1$ & $40$ & $18$ \\
$2$ & $40$ & $16$ \\
$3$ & $40$ & $20$ \\
\hline
\end{tabular}
\end{center}

If the town has $2{,}000$ residents, which is the best prediction for the number who prefer tea?
\end{questionText}
\choiceA{$800$}
\choiceB{$900$}
\choiceC{$1{,}000$}
\choiceD{$1{,}200$}
\correctAnswer{B}
\explanation{Find the average proportion: $\frac{18+16+20}{120} = \frac{54}{120} = 0.45$. Multiply: $0.45 \times 2{,}000 = 900$. About $900$ residents are predicted to prefer tea.}
\end{practiceQuestion}

\begin{practiceQuestion}{8-1-q25}{gsa}
\begin{questionText}
The table below shows the results of a survey asking $50$ randomly selected students about their favorite sport.

\begin{center}
\begin{tabular}{|l|c|}
\hline
\textbf{Sport} & \textbf{Students} \\
\hline
Soccer & $18$ \\
Basketball & $12$ \\
Swimming & $10$ \\
Tennis & $6$ \\
Other & $4$ \\
\hline
\end{tabular}
\end{center}

If the school has $800$ students, predict how many students prefer soccer.
\end{questionText}
\correctAnswer{$288$}
\explanation{The sample proportion for soccer is $\frac{18}{50} = 0.36$. Multiply by the population: $0.36 \times 800 = 288$. About $288$ students are predicted to prefer soccer.}
\end{practiceQuestion}

\begin{practiceQuestion}{8-1-q26}{gsa}
\begin{questionText}
Look at the two sampling scenarios described in the diagram below.

\begin{center}
\begin{tikzpicture}[scale=0.8]
% Scenario 1
\node[draw, rounded corners, fill=funBlue!15, minimum width=5cm, minimum height=1.2cm, align=center] at (0,2) {Scenario 1: Survey $30$ students\\waiting in the lunch line};
% Scenario 2
\node[draw, rounded corners, fill=funGreen!15, minimum width=5cm, minimum height=1.2cm, align=center] at (0,0) {Scenario 2: Randomly select $30$ students\\from the full student roster};
% Arrow labels
\node[right] at (3,2) {\textit{Convenience}};
\node[right] at (3,0) {\textit{Random}};
\end{tikzpicture}
\end{center}

Which scenario is more likely to produce a representative sample, and why?
\end{questionText}
\correctAnswer{Scenario 2, because every student has an equal chance of being selected.}
\explanation{Random sampling from the full roster gives every student an equal chance of selection, producing a representative sample. The lunch-line sample is convenience-based and may miss students who bring lunch or eat elsewhere.}
\end{practiceQuestion}

\begin{practiceQuestion}{8-1-q27}{gsa}
\begin{questionText}
A company surveyed $60$ of its $1{,}500$ employees about commute time. The results are shown below.

\begin{center}
\begin{tabular}{|l|c|}
\hline
\textbf{Commute Time} & \textbf{Number of Employees} \\
\hline
Less than $15$ min & $12$ \\
$15$ to $30$ min & $24$ \\
$31$ to $45$ min & $15$ \\
More than $45$ min & $9$ \\
\hline
\end{tabular}
\end{center}

Based on the sample, predict how many of the $1{,}500$ employees have a commute of $15$ to $30$ minutes.
\end{questionText}
\correctAnswer{$600$}
\explanation{The sample proportion is $\frac{24}{60} = 0.4$. Multiply by the population: $0.4 \times 1{,}500 = 600$. About $600$ employees are predicted to have a $15$- to $30$-minute commute.}
\end{practiceQuestion}
